\chapter{Simulation Uncertainty and Validity Diagnostics}
\label{app:diagnostics}

This appendix explains the rejection-rate denominators, uncertainty intervals
and validity checks used in Chapters~\ref{chap:design} and~\ref{chap:results}.

\section{Rejection rates and simulation uncertainty}

For each configuration, 20,000 table pairs are sampled. The plotted rate is
the number of valid rejections divided by 20,000. The method-valid conditional rate divides
valid rejections by that method's valid outputs; it is undefined if there are
none. The common-valid rate restricts to outputs where both methods are
valid. A paired difference is based on counts of the four joint rejection
outcomes for the same table pairs. These distinctions allow a reader to
check whether a curve difference is primarily a changed reference or lost
validity.

For the paired comparison, let \(D_s\) equal the Wald valid-rejection
indicator minus the Expanded Welch valid-rejection indicator for sampled
table pair \(s\), matching the sign used in
Table~\ref{tab:null-paired-examples}. Its average is
\(\bar D=(1/20{,}000)\sum_{s=1}^{20{,}000}D_s\), with approximate 95\% interval
\begin{equation}
  \bar D\ \pm\ 1.96
  \left\{
    \frac{1}{20{,}000(19{,}999)}
    \sum_{s=1}^{20{,}000}(D_s-\bar D)^2
  \right\}^{1/2}.
\end{equation}
This calculation preserves the within-replicate pairing.

The individual method intervals use Wilson's binomial score
construction \citep{wilson1927}.

\section{Validity and sparse-sample diagnostics}
\label{sec:validity-rules}

Both tests require matching table shapes with at least two rows and columns,
finite nonnegative integer counts, and more than one observation in each
table. Normal Wald requires a finite corrected MI difference, a positive
finite standard error, \(\Vhat(P)+\Vhat(Q)>10^{-14}\), and a finite p-value.
The variance threshold treats values extremely close to zero as numerically
unusable.

Expanded Welch adds component-level conditions to those shared requirements.
It requires positive finite degrees of
freedom for both populations and their combination. If one component's
degrees of freedom cannot be calculated, the expanded test is invalid even
when the direct combined expression in
Appendix~\ref{app:derivation-details} has a value. Failed calculations
contribute no rejection to the unconditional rate.

For any \(P\) population with sample size \(n_P\), its minimum expected
cell count is \(\min_{ij}n_Pp_{ij}\); define \(Q\)'s analogously. No
``expected minimum'' of random observed counts is calculated. The
sparsity diagnostics include empty observed rows and columns and the fraction
of zero-count cells. Such values describe why a statistic might fail, but
do not constitute a universal validity threshold.

The thesis does not treat a descriptive count of regimes closer to 0.05 as
a formal paired test over regimes. Wilson intervals quantify binomial Monte
Carlo variation at each fixed point, and paired normal intervals quantify a
specified within-configuration method difference. Neither interval family is
a simultaneous assertion across all configurations.

Each configuration also answers a distinct experimental question. Tables sharing a
population but changing \(n\) or \(|\Delta|\) answer different fixed-design
questions and should not be collapsed into one mean rate.

\section{Independent-pilot standard-error check}

The diagnostic in Section~\ref{sec:mechanism-results} uses two independent
sets of 20,000 sampled table pairs from the stated null populations. The
pilot set gives the sample standard deviation of the corrected MI
differences. The evaluation set is tested twice: once with each pair's
plug-in standard error, and once with that fixed pilot standard deviation.
Both calculations use the standard normal distribution and the two-sided
5\% threshold. The comparison of variance estimates divides the evaluation
set's average squared plug-in standard error by the squared pilot standard
deviation. This check retains the population pair and estimated differences
while changing how they are standardised.
